\documentclass{beamer}
\usefonttheme{professionalfonts}% Now beamer didn't modify the math fonts
\usepackage{geometry}
\input{macros.tex}
% \usepackage{tikz-cd}
\beamertemplatenavigationsymbolsempty
\setbeamertemplate{footline}[page number]{}
\setbeamerfont{footline}{size=\footnotesize}
\usepackage{tcolorbox}
\usepackage{biblatex}
% \usepackage{tikzcd}
\usetheme{}
\usepackage{tikz}

\definecolor{asparagus}{rgb}{0.53, 0.66, 0.42}


\title[]{Clase tutorial 10: Repaso de probabilidad y estadística.}
% \subtitle{Workshop de LOREL 2024.}
\date{22 de Octubre de 2024.}
\author[]{ Investigación operativa.}
\institute{Universidad de San Andrés}
% \titlegraphic{\hfill\includegraphics[height=1.5cm]{logo.pdf}}

\begin{document}
\maketitle

\begin{frame}{Objetivos de la clase de hoy.}
  La idea de la clase de hoy es ver los siguientes temas:
  \begin{itemize}
    \item Seguir repasando temas de probabilidad y estadística.
  \end{itemize}
\end{frame}

\begin{frame}[fragile]{Problema 1.}
  La vida util en días de las lamparitas marca Osram tiene una distribución exponencial de media 100 y la de las lamparidas marca Philips tiene una distribución exponencial de media 150. 
  Si nos regalaron una lamparita (la eligieron al azar sin mirar la marca, pero se sabe que es Osram o Philips) y la estamos usando hace a 200 días sin que se queme, ¿cuál es la probabilidad de que sea Philips?
\end{frame}


\begin{frame}[fragile]{Problema 2.}
  El modelo binomial para el precio de una acción supone que si el precio presente de la acción es $S_0$, entonces luego de un periodo su valor será $S_1 = uS_0$ con probabilidad $p$ y $S_1 = dS_0$ con probabilidad $1 - p$.
Asumimos que los movimientos del precio en cada periodo son independientes. Asumimos para simplificar que $d = \frac{1}{u}$ y $p = 0.5$, $u = 1.012$.

Decir como está distribuida $S_n$, la variable aleatoria que nos da el precio de la acción luego de $n$ periodos.

Para $n = 1000$, calcular la probabilidad de que el precio de la acción sea un $30\%$ mayor al inicial $S_0$. 
Hacerlo también aproximando con una normal.
\end{frame}

\begin{frame}[fragile]{Problema 3.}
  $6$ amigos se quieren juntar en una pizzería y hacen una reserva para las $9:00 \text{ pm}$. Si todos llegan $10$ minutos tarde pierden la reserva. Supongamos que el tiempo de llegada de los amigos es independiente, centrado en las $9:00 \text{ pm}$ y esta distribuido como $\mathcal{U}(-5,25)$. 
  
  ¿Cuál es la probabilidad de que pierdan la reserva?


\end{frame}


\begin{frame}[fragile]{Tarea.}
  
La cantidad de insectos que arriban a la mesa de un asado responde
a una distribución Poisson de media $\lambda = 60$ por hora. Cada insecto
puede ser una mosca con probabilidad $p = 2/3$. La cantidad de insectos
y la naturaleza de cada uno de ellos son independientes.

\begin{enumerate}[a)]
  \item Dado que arribaron 60 insectos, ¿cual es la distribución de la
  cantidad de moscas en la mesa?
  
  \item Hallar la probabilidad de que arriben 50 insectos y 35 de ellos
  sean moscas.
\end{enumerate}
\end{frame}

\end{document}
